Calcul arithmétique, pour une inflation de 2 \% et les trois régimes français de 2026 : 17,2 \% de prélèvements sociaux pour l'assurance vie dans l'abattement, 31,4 \% pour le prélèvement forfaitaire unique, 47,2 \% pour un revenu foncier taxé au taux marginal de 30 \%.\par Note : l'impôt frappe les intérêts nominaux, dont une partie ne fait que compenser la hausse des prix ; il est donc dû même quand l'épargnant ne s'est pas enrichi. Le rendement nominal au-dessous duquel épargner appauvrit vaut l'inflation divisée par un moins le taux. La partie colorée de chaque barre est le rendement réel minimal que le placement doit dégager : 0,92 \% au prélèvement forfaitaire unique, contre 0 si l'impôt portait sur le rendement réel.
